% Hafta 11 -- alıştırma gövdesi (ders notu ve ayrı alıştırma kağıdı tarafından \input edilir)
\section{Çözümlü Alıştırmalar}
Her soruyu önce kendiniz deneyin, sonra çözümü okuyun.

\begin{alistirma}
\(12a^2b-18ab^2\) ifadesini çarpanlarına ayırınız.
\end{alistirma}
\begin{cozum}
EBOB\((12,18)=6\); her iki terimde de \(a\) ve \(b\) var \(\Rightarrow\) ortak çarpan \(6ab\).
\[12a^2b-18ab^2=6ab\,(2a-3b)\]
Kontrol: \(6ab\cdot 2a=12a^2b\), \ \(6ab\cdot 3b=18ab^2\) \checkmark
\end{cozum}

\begin{alistirma}
\(x^2-y^2+3x-3y\) ifadesini çarpanlarına ayırınız.
\end{alistirma}
\begin{cozum}
İlk iki terim iki kare farkı, son iki terimde ortak çarpan \(3\):
\[x^2-y^2+3x-3y=(x-y)(x+y)+3(x-y)\]
Ortak parantez \((x-y)\) dışarı alınır:
\[=(x-y)(x+y+3)\]
\end{cozum}

\begin{alistirma}
\((x+1)^2-25\) ifadesini çarpanlarına ayırınız.
\end{alistirma}
\begin{cozum}
\(a=x+1\), \(b=5\) olmak üzere iki kare farkı:
\[(x+1)^2-5^2=(x+1-5)(x+1+5)=(x-4)(x+6)\]
Kontrol: \((x-4)(x+6)=x^2+2x-24\) ve \((x+1)^2-25=x^2+2x+1-25=x^2+2x-24\) \checkmark
\end{cozum}

\begin{alistirma}
\(a-b=5\) ve \(a^2+b^2=37\) olduğuna göre \(a\cdot b\) kaçtır?
\end{alistirma}
\begin{cozum}
\((a-b)^2=a^2-2ab+b^2\) özdeşliğinde verilenler yerine yazılır:
\[5^2=37-2ab\ \Rightarrow\ 25=37-2ab\ \Rightarrow\ 2ab=12\ \Rightarrow\ ab=6\]
(Kontrol: \(a=6,\ b=1\): \(a-b=5\), \(36+1=37\), \(ab=6\).)
\end{cozum}

\begin{alistirma}
\(x^2+kx+49\) bir tam kare ve \(k<0\) olduğuna göre \(k\) kaçtır?
\end{alistirma}
\begin{cozum}
\(x^2=(x)^2\), \(49=7^2\). Tam kare ifade \((x+7)^2\) veya \((x-7)^2\) olabilir.\\
\(k<0\) olduğundan \((x-7)^2=x^2-14x+49\) \(\Rightarrow\) \(k=-14\).
\end{cozum}

\begin{alistirma}
\(4x^2-4x-3\) ifadesini çarpanlarına ayırınız.
\end{alistirma}
\begin{cozum}
\begin{minipage}[t]{0.2\linewidth}
\begin{tabular}{ccc}
\(2x\) & & \(-3\)\\
\(2x\) & & \(1\)
\end{tabular}
\end{minipage}\hfill
\begin{minipage}[t]{0.62\linewidth}
\(4x^2=2x\cdot 2x\), \ \(-3=(-3)\cdot 1\).\\
Çapraz: \(2x\cdot 1+2x\cdot(-3)=2x-6x=-4x\) \checkmark
\[4x^2-4x-3=(2x-3)(2x+1)\]
\end{minipage}
\end{cozum}

\begin{alistirma}
\(a+b=6\) ve \(a\cdot b=5\) olduğuna göre \(a^3+b^3\) kaçtır?
\end{alistirma}
\begin{cozum}
\[a^3+b^3=(a+b)^3-3ab(a+b)=6^3-3\cdot 5\cdot 6=216-90=126\]
İkinci yol: \(a^2+b^2=36-10=26\); \ \(a^3+b^3=(a+b)(a^2-ab+b^2)=6\cdot(26-5)=126\).\\
(Kontrol: \(a=1,\ b=5\): \(1+125=126\).)
\end{cozum}

\begin{alistirma}
\(\dfrac{x^3+8}{x^2-4}\) ifadesini sadeleştiriniz ve tanım koşulunu yazınız.
\end{alistirma}
\begin{cozum}
Pay küp toplamı, payda iki kare farkı:
\[\frac{x^3+8}{x^2-4}=\frac{(x+2)(x^2-2x+4)}{(x-2)(x+2)}=\frac{x^2-2x+4}{x-2}\]
Tanım koşulu ilk paydadan: \(x^2-4\neq 0\) \(\Rightarrow\) \(x\neq 2\) ve \(x\neq -2\).
\end{cozum}

\kendinizoncesi
\section{Alıştırmalar (Kendiniz Çözün)}

\begin{enumerate}[label=\textbf{\arabic*.},leftmargin=*,itemsep=6pt]
\item \(7x^2-14x\) ifadesini çarpanlarına ayırınız.
\item \(xy+2x+3y+6\) ifadesini çarpanlarına ayırınız.
\item \(4x^2-81\) ifadesini çarpanlarına ayırınız.
\item \(75^2-25^2\) işleminin sonucu kaçtır?
\item \(x^4-81\) ifadesini çarpanlarına ayırınız.
\item \((2a+5)^2\) açılımını yazınız.
\item \(a+b=9\) ve \(a\cdot b=14\) olduğuna göre \(a^2+b^2\) kaçtır?
\item \(x+\dfrac{1}{x}=6\) olduğuna göre \(x^2+\dfrac{1}{x^2}\) kaçtır?
\item \(x^2-10x+m\) bir tam kare ise \(m\) kaçtır?
\item \(x^2+5x-14\) ifadesini çarpanlarına ayırınız.
\item \(2x^2-5x-3\) ifadesini çarpanlarına ayırınız.
\item \(x^2-6x+8=0\) denkleminin çözüm kümesini bulunuz.
\item \(x^3-27\) ifadesini çarpanlarına ayırınız.
\item \(a-b=1\) ve \(a\cdot b=6\) olduğuna göre \(a^3-b^3\) kaçtır?
\item \(\dfrac{x^2+2x-8}{x^2-4}\) ifadesini sadeleştiriniz ve tanım koşulunu yazınız.
\item \(x+\dfrac{1}{x}=4\) olduğuna göre \(x^3+\dfrac{1}{x^3}\) kaçtır?
\end{enumerate}

\needspace{10\baselineskip}
\subsection*{Cevap Anahtarı}
{\small
\begin{multicols}{2}
\begin{enumerate}[label=\textbf{\arabic*.},leftmargin=*,itemsep=1pt]
\item Ortak çarpan \(7x\): \(7x(x-2)\).
\item \(x(y+2)+3(y+2)=(y+2)(x+3)\).
\item \((2x)^2-9^2=(2x-9)(2x+9)\).
\item \((75-25)(75+25)=50\cdot 100=5000\).
\item \((x^2-9)(x^2+9)=(x-3)(x+3)(x^2+9)\).
\item \(4a^2+20a+25\).
\item \(9^2-2\cdot 14=81-28=53\).
\item \(6^2-2=34\).
\item \((x-5)^2\Rightarrow m=25\).
\item Çarpım \(-14\), toplam \(5\): \((x+7)(x-2)\).
\item Çapraz: \((2x+1)(x-3)\).
\item \((x-2)(x-4)=0\Rightarrow\{2,\,4\}\).
\item \((x-3)(x^2+3x+9)\).
\item \(1^3+3\cdot 6\cdot 1=19\).
\item \(\dfrac{(x+4)(x-2)}{(x-2)(x+2)}=\dfrac{x+4}{x+2}\); \ \(x\neq\pm 2\).
\item \(4^3-3\cdot 4=64-12=52\).
\end{enumerate}
\end{multicols}}
